\cotefile{105}
\scanpage{37}
% archive sheet number: 36

\handwritten{UTDE}\\
\handwritten{UTIL}\\
\handwritten{UTSD}

% printed form
{\tiny\ttfamily
\noindent\begin{tabular}{@{}l@{\ }llllllll@{\ }l@{}}
&\multicolumn{8}{l}{llpd 'cnu1.modulef.utda'} &\\
& AFMDES & CLOSE & \tikz[baseline=(t.base)]\node[draw,ellipse,inner sep=0.3pt](t){\tiny\ttfamily CUTTER}; & \tikz[baseline=(d.base)]{\node[inner sep=0pt](d){\tiny\ttfamily DYNAM};\handwritten{\draw[thin] (d.west) to[bend right=55] (d.south);}} & ECRDIR & ECR1 & FLIPFLOP & INCANO & \handwritten{21}\\
\handwritten{\unsure{1|I}} & INCAPA & INDUNI & IVIRT & LAYOUT & LECDIR & LEC1 & LISTPDS & MDES &\\
\handwritten{B} & MESSAGE & NUMALP & OPEN & TRANSLAT & UPDATE &&&&\\
\handwritten{$\Pi$} & \multicolumn{8}{l}{llpd 'cnu1.modulef.utdas'} &\\
& AFMDES & CLOSE & CUTTER & DYNAM & ECRDIR & FLIPFLOP & INCANO & INDUNI &\\
& IVIRT & LAYOUT & LECDIR & LISTPDS & MDES & MESSAGE & NUMALP & OPEN &\\
& TRANSLAT & UPDATE &&&&&&&\\ \hline
\end{tabular}}

{\tiny\ttfamily
\noindent\begin{tabular}{@{}llllllll@{\ }l@{}}
\multicolumn{8}{l}{llpd 'cnu1.modulef.utde'} &\\
AFMDES & CLOSE & CUTTER & DEFDIR & DYNAM & ECPDIR & ECP1 & ECRDIR & \handwritten{51}\\
ECRP0 & ECR1 & ECTDIR & ERLIBR & FERDIR & FLIPFLOP & INCANO & INCAPA &\\
INDUNI & INIDIR & INITI & IVIRT & LAYOUT & LECDIR & LECTAB & LECTUD &\\
LECTUI & LECTUK & LECTUT & LECTU1 & LECTU2 & LECTU3 & LECTU4 & LECTU5 &\\
LECT10 & LECVAL & LEC1 & LEPDIR & LEP1 & LETDIR & LISTPDS & MDES &\\
MESS & MESSAGE & NBLIGN & NBMACH & NBMOTS & NUMALP & OPEN & OUVDIR &\\
POLIBR & TRANSLAT & UPDATE &&&&&&\\
\multicolumn{8}{l}{llpd 'cnu1.modulef.utdes'} &\\
ECTDIR & ERLIBR & INCANO & INITI & LECTAB & LECTUD & LECTUI & LECTUK & \handwritten{20}\\
LECTUT & LECTU1 & LECTU2 & LECTU3 & LECTU4 & LECTU5 & LECT10 & LECVAL &\\
LETDIR & MESS & NBLIGN & POLIBR &&&&&\\ \hline
\end{tabular}}

{\tiny\ttfamily
\noindent\begin{tabular}{@{}l@{\ }llllllll@{\ }l@{}}
& \multicolumn{8}{l}{llpd 'cnu1.modulef.utdf'} &\\
\handwritten{1} & DEFDIR & ECPDIR & ECP1 & ECRP0 & FERDIR & INIDIR & LEPDIR & LEP1 & \handwritten{11}\\
\handwritten{B} & NBMACH & NBMOTS & OUVDIR &&&&&&\\
& \multicolumn{8}{l}{llpd 'cnu1.modulef.utdfs'} &\\
\handwritten{$\Omega$} & DEFDIR & ECPDIR & ECRP0 & FERDIR & INIDIR & LEPDIR & NBMACH & NBMOTS &\\
& OUVDIR &&&&&&&&\\ \hline
\end{tabular}}

{\tiny\ttfamily
\noindent\begin{tabular}{@{}llllllll@{}}
\multicolumn{8}{l}{llpd 'cnu1.modulef.utdm'}\\
\multicolumn{8}{l}{TSC0185I ANGIVET BIBLIOTEK (PDS) HAR INGA MEDLEMMAR}\\
\multicolumn{8}{l}{llpd 'cnu1.modulef.utdms'}\\
\$RETURN & \$SAVE & \$STOP & \pounds COMDIR & \pounds DBDAM & \pounds DBSAM & \pounds DES & CALL\\
CAMLST & CHECK & CLOSE & DCBD & FREEMAIN & GETMAIN & IECSDSL1 & IECTDECB\\
IEFASCTB & IEFASIOT & IEFJFCBN & IEFTIOT1 & IEZBITS & IEZDEB & IEZJSCB & IHACDE
\end{tabular}}

\scanpage{38}
% archive sheet number: 37

\mspage{15}
\uline{Corollary} Assume \deleted{$\Phi_0$ is \unsure{sm}}

a) $\Phi_0$ is small, source and target of any $g\in\Phi_0$ are accessible

b) $M$ stable by filtering direct limits, \deleted{\unsure{and are|and is|and} by \unsure{\illegible{2}|amalgamated}} \deleted{\unsure{pres by amlg|push by analogy}}

\textcircled{c} \deleted{\unsure{having amalgamated|having an amalgamated|Any amalgamated}} \inserted{For any} diagram

\diagram{\begin{tikzcd}[ampersand replacement=\&]
 A \arrow[r,no head,"\unsure{v|u}"] \arrow[d,"\alpha"'] \& X\\
 B \&
\end{tikzcd}}

with $\alpha\in\Phi_{\deleted{\unsure{0}}}$, the amalgamated sum exists

\diagram{\begin{tikzcd}[ampersand replacement=\&]
 A \arrow[r,"v"] \arrow[d,"\alpha"'] \& X \arrow[d,"\alpha'"]\\
 B \arrow[r,"u'"'] \& X'
\end{tikzcd}}

\margin{Also c') $\Phi$ stable under compositions}

and $\alpha'\in\Phi$

\textcircled{d} $\Phi$ ``stable by \inserted{\uline{filtering}} direct limits''

\margin{It is enough to have it wrt. ordered sets $I$ where Sup. of any two elements exists}

Then $(\Phi,\Psi)$ satisfy the factorisation condition, and \deleted{\unsure{even|more}} even, there exists a functor $\Sigma$.

\deleted{Ex.}

\diagram{\begin{tikzcd}[ampersand replacement=\&]
 \underline{\mathrm{Fl}}(M)\arrow[rr,no head,"\Sigma"]\arrow[dr,"t"] \& \& \underline{\mathrm{Fl}}(M)\arrow[dl,"t"']\\
 \& M \&
\end{tikzcd}},

and a functorial morphism \deleted{\unsure{$\xi$}} in $[\underline{\mathrm{Cat}}]_{/M}$

\margin{possibly ``large'' categories over $M$}

\[
\deleted{\unsure{\phi\longrightarrow}}\ \mathrm{id}_{\underline{\mathrm{Fl}}(M)}\longrightarrow\Sigma.
\]

\diagram{\begin{tikzcd}[ampersand replacement=\&]
 X \arrow[r,"i(f)"]\arrow[d] \& \overline X(f)\arrow[dl,"p(f)"]\\
 Y \&
\end{tikzcd}}

with $i(f)\in\Phi,\ p(f)\in\Psi$

\uline{Example} Assume that $M$ is stable by amalgamated sums. Take $\Phi=\overline{\Phi_0}=\Psi^*$. Then c) \deleted{and d)} is trivially satisfied

What about d)? OK \inserted{for \uline{countable} \uline{limits}}. To get \deleted{\unsure{extensions of|extension of}} \inserted{the} dotted $\varphi$,

\diagram{\begin{tikzcd}[ampersand replacement=\&]
 A_{i_0}\arrow[r,no head]\arrow[d,"u_0"'] \& A_i\arrow[r] \& B=\underset{\longrightarrow}{\lim}\, A_i\arrow[d,"v"]\arrow[dll,dashed,"\varphi"']\\
 X\arrow[rr,"f","{\rotatebox{90}{$\in$}\,\Psi}"'] \& \&Y
\end{tikzcd}}

we make a Zorn induction on pairs $(\inserted{\overset{\text{\unsure{filtering positive|filtering quite}}}{J}}\in I,\ u_J\colon \overbrace{\underset{i\in J}{\underset{\longrightarrow}{\lim}}\ A_i}^{\inserted{A_J}}\to X)$ which make the diagrams

\diagram{\begin{tikzcd}[ampersand replacement=\&]
 A_{i_0}\arrow[r,no head]\arrow[rr,bend left=13,no head]\arrow[d,"u_0"'] \& A_J\arrow[r]\arrow[dl,"u_J"'] \& B\arrow[d]\\
 X\arrow[rr,"f"] \& \&Y
\end{tikzcd}}

commute, with evident order relation. \deleted{This is clearly} inductive.

\[
 \mathrm{Hom}(B,X)\longrightarrow
 \underset{\underset{\leftarrow}{\lim}\,\mathrm{Hom}(g_i,f)}{\underset{\scriptscriptstyle\shortparallel}{\mathrm{Hom}(g,f)}}\quad\text{surjection}
\]

\scanpage{39}
% archive sheet number: 38

% printed form
{\tiny\ttfamily
\noindent\begin{tabular}{@{}llllllll@{}}
IHALLE & IHAPSA & IHARB & IHASPQE & IHAXTLST & IKJRB & IKJTCB & OBTAIN\\
OPEN & RDJFCB & READ & SYNADAF & SYNADRLS & WRITE
\end{tabular}}

{\tiny\ttfamily
\noindent\begin{tabular}{@{}llllllll@{}}
\multicolumn{8}{l}{llpd 'cnu1.modulef.util'}\\
ABEGCD & ABEGCR & AB0D & AB0R & AB1D & AB1R & ADRESS & AFFCOM\\
AFFKOM & ALEAD & ALEAR & ASMPVD & ATAD & ATAR & ATB0D & ATB0R\\
ATB1D & ATB1R & ATMPBD & ATMPVD & AZEROD & AZEROR & BASYMD & CADRE\\
CHOLOC & CL2TBD & CL2TBR & CL2VED & CL2VER & COMPRE & CROLOC & DECAL\\
DECODA & DET3D & DGELG & DILATE & FOARET & ICHIF & IEGAL & IMATAB\\
IMPRIA & IMPTAB & INVCRO & INVLOC & ISIGNO & ISOMME & MAJDCB & MAXTAM\\
MNEMO & MUPLIN & MUPOB2 & MXPOIR & MXPO1R & NBARET & NBDIFF & NBFACE\\
NBMAXI & NBMINI & NBSOME & NEWS & NOM & NOMAFF & NOMLON & NUMALP\\
PA & PB & PLMASD & PLONAD & PN2DDE & PN2DIN & PN2DPR & PN2DVA\\
PN3DDE & PN3DIN & PN3DPR & PN3DVA & PROJEC & PROLIG & PROSCD & PROSCR\\
PU & PV & PW & PY & P1 & Q1 & RCHLOC & READRE\\
RENOMM & SECANT & SGNHEX & SGNQUA & SGNTET & SGNTRI & SOMABD & SOMABR\\
SOMMEM & TABAND & TABANR & TABARD & TABARR & TABASD & TABASR & TABAUD\\
TABAUR & TABAZD & TABAZR & TABA0D & TABA0R & TABA1D & TABA1R & TABA2D\\
TABA2R & TABA3D & TABA3R & TABA4D & TABA4R & TABA5D & TABA5R & TABA6D\\
TABA6R & TABA7D & TABA7R & TABA8D & TABA8R & TABA9D & TABA9R & TABSUD\\
TABSUR & TABSZD & TABSZR & TAB0D & TAB0R & TAB1D & TAB1R & TRCROI\\
TRFLEC & TROUVE & TRTATA & TRTAVE & TUER & TYMAEL & TYPCLA & TYPGEL\\
TYPMOT & VERIF
\end{tabular}}

% slanted
\margin{\handwritten{\unsure{utilitaires}\\ fortran\\ 146}}

{\tiny\ttfamily
\noindent\begin{tabular}{@{}llllllll@{}}
\multicolumn{8}{l}{llpd 'cnu1.modulef.utils'}\\
ABEGCD & ABEGCR & AB0D & AB0R & AB1D & AB1R & ADRESS & AFFCOM\\
AFFKOM & ALEAD & ALEAR & ASMPVD & ATAD & ATAR & ATB0D & ATB0R\\
ATB1D & ATB1R & ATMPBD & ATMPVD & AZEROD & AZEROR & BASYMD & CADRE\\
CHOLOC & CL2TBD & CL2TBR & CL2VED & CL2VER & COMPRE & CROLOC & DECAL\\
DECODA & DET3D & DGELG & DILATE & FOARET & ICHIF & IEGAL & IMATAB\\
IMPRIA & IMPTAB & INVCRO & INVLOC & ISIGNO & ISOMME & MAJDCB & MAXTAM\\
MNEMO & MUPLIN & MUPOB2 & MXPOIR & MXPO1R & NBARET & NBDIFF & NBFACE\\
NBMAXI & NBMINI & NBSOME & NEWS & NOM & NOMAFF & NOMLON & NUMALP\\
PA & PB & PLMASD & PLONAD & PN2DDE & PN2DIN & PN2DPR & PN2DVA\\
PN3DDE & PN3DIN & PN3DPR & PN3DVA & PROJEC & PROLIG & PROSCD & PROSCR\\
PU & PV & PW & PY & P1 & Q1 & RCHLOC & READRE\\
RENOMM & SECANT & SGNHEX & SGNQUA & SGNTET & SGNTRI & SOMABD & SOMABR\\
SOMMEM & TABAND & TABANR & TABARD & TABARR & TABASD & TABASR & TABAUD\\
TABAUR & TABAZD & TABAZR & TABA0D & TABA0R & TABA1D & TABA1R & TABA2D\\
TABA2R & TABA3D & TABA3R & TABA4D & TABA4R & TABA5D & TABA5R & TABA6D\\
TABA6R & TABA7D & TABA7R & TABA8D & TABA8R & TABA9D & TABA9R & TABSUD\\
TABSUR & TABSZD & TABSZR & TAB0D & TAB0R & TAB1D & TAB1R & TRCROI\\
TRFLEC & TROUVE & TRTATA & TRTAVE & TUER & TYMAEL & TYPCLA & TYPGEL\\
TYPMOT & VERIF\\ \hline
\end{tabular}}

{\tiny\ttfamily
\noindent\begin{tabular}{@{}llllllll@{}}
\multicolumn{8}{l}{llpd 'cnu1.modulef.utsd'}\\
AELESU & AJASS & CNDLN & COPISD & COSDB & COTASD & CTYTAE & CTYTA1\\
ECAMAT & ECATRI & ECB & ECBDCL & ECCOOR & ECFORC & ECHYMA & ECMAIL\\
ECMILI & ECMPLR & ECMUA & ECNOPO & ECRIA1 & ECRIT1 & ECRTAE & ECTAE\\
ECTASD & ECTYNO & ECVISU & ECVVPR & ESTASF & ESVVPR & IAIJ1 & IMAMAT\\
IMAMA1 & IMAMA2 & IMATRI & IMB & IMBDCL & IMB1 & IMCOOR & IMCOO1\\
IMFORC & IMHYMA & IMMAIL & IMMAI1 & IMMILI & IMMPLR & IMMUA & IMNOPO\\
IMTADE & IMTAE & IMTASD & IMTA01 & IMTYNO & IMVISU & IMVVPR & IM1BDC\\
INICSD & INITMP & INTAB0 & ISNOP5 & LEAMAT & LEAMMS & LEATRI & LEB\\
LEBDCL & LECOOR & LEFORC & LEHYMA & LEMAIL & LEMEFI & LEMILI & LEMPLR\\
LEMUA & LENOPO & LETAE & LETASD & LETYNO & LEVISU & LEVVPR & LRNDL1\\
LSVVPR & MONOPO & MONOP5 & MONOP6 & NCMUNS & NDL & NIVLIB & NUPAG\\
NUPAGE & PRINTB & PRIN1B & RACRNO & REAMAT & REATRI & REB & REBDCL\\
RECOOR & REFORC & REHYMA & REMAIL & REMEFI & REMILI & REMPLR & REMUA\\
RENOPO & RESOUR & RETAE & RETASD & RETYNO & REVISU & REVVPR & RSTSDE\\
SAUCSD & SAUSDS & SAUTP & SDBMSC & STAMD & STAMR & TRAMAD & TRAMAR\\
TUERSD & TUTASD & VERNN & VERTNN
\end{tabular}}

% slanted
\margin{\handwritten{fortran\\ 124}}

{\tiny\ttfamily
\noindent\begin{tabular}{@{}llllllll@{}}
\multicolumn{8}{l}{llpd 'cnu1.modulef.utsds'}\\
AELESU & AJASS & CNDLN & COPISD & COSDB & COTASD & CTYTAE & CTYTA1\\
ECAMAT & ECATRI & ECB & ECBDCL & ECCOOR & ECFORC & ECHYMA & ECMAIL\\
ECMILI & ECMPLR & ECMUA & ECNOPO & ECRIA1 & ECRIT1 & ECRTAE & ECTAE\\
ECTASD & ECTYNO & ECVISU & ECVVPR & ESTASF & ESVVPR & IAIJ1 & IMAMAT\\
IMAMA1 & IMAMA2 & IMATRI & IMB & IMBDCL & IMB1 & IMCOOR & IMCOO1\\
IMFORC & IMHYMA & IMMAIL & IMMAI1 & IMMILI & IMMPLR & IMMUA & IMNOPO
\end{tabular}}

\scanpage{40}
% archive sheet number: 39

\mspage{16}
\uline{Lemma} Let $\Phi\subset\mathrm{Fl}(M)$ be left $Q$-saturated (i.e. $\Phi=\Psi^*$, \inserted{for some $\Psi$})

then \deleted{$\Phi$ is stable by composition} \unsure{(}

a) $\Phi$ contains iso

b) $\Phi$ stable \deleted{by} \inserted{under} composition and cobase changes

c) $\Phi$ stable \deleted{by direct} \deleted{\unsure{order}} \inserted{under} arbitrary \inserted{(including infinite)} amalgamated sums

% slanted
\margin{\uline{NB} $\overset{a,b,d}{\Longrightarrow}$ c)}

d) $\Phi$ stable \deleted{by direct systems indexed by any well ordered set -- i.e. if we have any direct system} \deleted{\unsure{under well ordered|directed ordered set}} \inserted{under ordinal direct systems, i.e.} if we have any well ordered set $I$ and $(A_i)_{i\in I}$ direct system, with $g_{ij}\colon A_i\to A_j$ in $\Phi$, \inserted{and $A_{i_0}=\underset{j<i_0}{\underset{\longrightarrow}{\lim}} A_j$ if $i_0\in I$ is a limit,} then \inserted{$\forall i_0\in I$} $A_{i_0}\xrightarrow{g_{i_0}}B=\underset{\longrightarrow}{\lim}\, A_i$ is in $\Phi$.

\deleted{\unsure{Then|The} only}

\uline{Thus we} get,

\uline{Thm} Let $M$ be a category, stable by small direct limits and (weakly) accessible. Then

\textcircled{1} \deleted{\unsure{a}} Let $\Phi_0^{\deleted{\unsure{*}}}\subset\mathrm{Fl}(M)$,$^*$ and $\Phi_0\subset\Phi\subset\overline{\Phi_0}\overset{\mathrm{dfn}}{=}\Psi^*$ (with $\Psi=\Phi_{0*}$). If $\Phi$ \deleted{stable \unsure{under|und|and}} satisfies a) b) c) d),$^*$ then the $(\Phi,\Psi)$ satisfies factorization condition. (Hence, $\Phi$ stable by direct factors, $\Phi=\overline{\Phi_0}$)

% slanted
\margin{We have to assume $\Phi_0$ small \deleted{\unsure{(?)}}.}

% slanted
\margin{We may ignore c) (which results from a) b) d)) \uline{or} assume only a, b', c, d, where b') is stability under cobase change (without assuming \deleted{\unsure{the}} stability under composition)}

\textcircled{2} \inserted{(same data $\Phi_0,\Phi$.)} In order that $\Phi$ \deleted{$=\overline{\Phi_0}$} \inserted{be saturated}, it is n.s. that $\Phi$ satisfy conditions a) b) c) d) and e) stability by direct factors. In this case, factorization holds for the $Q$-orthogonal pair $(\Phi,\Psi)$.

NB. In 2, \deleted{besides} there was the preliminary assumption of existence of $\Phi_0\subset\Phi$ s.th. $\Phi_0$ and $\Phi$ have the same left $Q$-saturation, with $\Phi_0$ \uuline{small}. Next step is to try to get rid of this.

\scanpage{41}
% archive sheet number: 40

% printed form
{\fontsize{4.1pt}{4.8pt}\selectfont
\begin{verbatim}
           00000000      00000000      00000000      00000000
         0000000000    0000000000    0000000000    0000000000
       00      0000  00      0000  00      0000  00      0000
      00     00 00  00     00 00  00     00 00  00     00 00
     00    00  00  00    00  00  00    00  00  00    00  00
    00   00   00  00   00   00  00   00   00  00   00   00
   00  00    00  00  00    00  00  00    00  00  00    00
  00 00     00  00 00     00  00 00     00  00 00     00
 0000      00  0000      00  0000      00  0000      00
000       00  000       00  000       00  000       00
0000000000    0000000000    0000000000    0000000000
00000000      00000000      00000000      00000000
\end{verbatim}
}

{\fontsize{4.1pt}{4.8pt}\selectfont
\begin{verbatim}
PPPPPPPPPPP   RRRRRRRRRRR   EEEEEEEEEEEE  DDDDDDDDD     OOOOOOOOOOOO  NN        NN  TTTTTTTTTTTT   AAAAAAAAAA
PPPPPPPPPPPP  RRRRRRRRRRRR  EEEEEEEEEEEE  DDDDDDDDDD    OOOOOOOOOOOO  NNN       NN  TTTTTTTTTTTT  AAAAAAAAAAAA
PP        PP  RR        RR  EE            DD       DD   OO        OO  NNNN      NN       TT       AA        AA
PP        PP  RR        RR  EE            DD        DD  OO        OO  NN NN     NN       TT       AA        AA
PP        PP  RR        RR  EE            DD        DD  OO        OO  NN  NN    NN       TT       AA        AA
PPPPPPPPPPPP  RRRRRRRRRRRR  EEEEEEEE      DD        DD  OO        OO  NN   NN   NN       TT       AAAAAAAAAAAA
PPPPPPPPPPP   RRRRRRRRRRR   EEEEEEEE      DD        DD  OO        OO  NN    NN  NN       TT       AAAAAAAAAAAA
PP            RR    RR      EE            DD        DD  OO        OO  NN     NN NN       TT       AA        AA
PP            RR     RR     EE            DD        DD  OO        OO  NN      NNNN       TT       AA        AA
PP            RR      RR    EE            DD       DD   OO        OO  NN       NNN       TT       AA        AA
PP            RR       RR   EEEEEEEEEEEE  DDDDDDDDDD    OOOOOOOOOOOO  NN        NN       TT       AA        AA
PP            RR        RR  EEEEEEEEEEEE  DDDDDDDDD     OOOOOOOOOOOO  NN         N       TT       AA        AA
\end{verbatim}
}

{\fontsize{4.1pt}{4.8pt}\selectfont
\begin{verbatim}
****A   END   JOB  637  PREDONTA                        ROOM       11.28.14 AM 22 APR 82  R6.PR1   SYS JES2  JOB  637   END   A****
****A   END   JOB  637  PREDONTA                        ROOM       11.28.14 AM 22 APR 82  R6.PR1   SYS JES2  JOB  637   END   A****
****A   END   JOB  637  PREDONTA                        ROOM       11.28.14 AM 22 APR 82  R6.PR1   SYS JES2  JOB  637   END   A****
****A   END   JOB  637  PREDONTA                        ROOM       11.28.14 AM 22 APR 82  R6.PR1   SYS JES2  JOB  637   END   A****
****A   END   JOB  637  PREDONTA                        ROOM       11.28.14 AM 22 APR 82  R6.PR1   SYS JES2  JOB  637   END   A****
****A   END   JOB  637  PREDONTA                        ROOM       11.28.14 AM 22 APR 82  R6.PR1   SYS JES2  JOB  637   END   A****
\end{verbatim}
}

\scanpage{43}
% archive sheet number: 42

\mspage{6}
\deleted{$q\in W\Longleftrightarrow q_{\overline U}\in W$}

\deleted{Or, on a}

\deleted{\diagram{\begin{tikzcd}[ampersand replacement=\&]
 U\arrow[r,no head,"i_0"]\arrow[d,no head,"q_U"'] \& \overline U\arrow[d,"q_{\overline U}"]\\
 U'\arrow[r,no head,"i'_0"'] \& \overline U'
\end{tikzcd}}}

\deleted{avec $i_0,i'_0\in W$, donc $q_{\overline U}\in W$ ssi $q_U\in W$. Donc $q\in W$ ssi $q_U\in W$, \unsure{q.e.d.|cqfd.}}

\uline{Bilan}. de W7, W(\deleted{\unsure{3}}7bis) si $W=W_{\underline D}$ ?

Se donner des $X$ ouverts $U,U'$ \deleted{avec $U\cup U'=X$}, revient au $=$ que de se donner $X\xrightarrow{\varphi}\Delta^{\unsure{1|'}}\times\Delta^1$, les fibres des quatre points \unsure{sts} resp. \deleted{\unsure{et|étant}},

\[
\left\{\begin{aligned}
 X_{00}&=U\cap U',\quad X_{01}=U\cap T',\quad X_{1,0}=T\cap U',\quad X_{1,1}=T\cap T'\\
 U&=\varphi^{-1}(\{0\}\times\Delta^1),\quad U'=\varphi^{-1}(\Delta^1\times\{0\})
\end{aligned}\right.
\]

\deleted{Donc} ($T'=X\smallsetminus U,\ T'=X\smallsetminus U'$). \deleted{Donc} $U\cup U'=X$ i.e. $T\cap T'=\phi$ équivaut à $X_{11}=\phi$, donc au fait qu'on a un foncteur $X\to\Psi=\Delta^1\times\Delta^1\smallsetminus\{1,1\}$

\diagram{\begin{tikzcd}[ampersand replacement=\&,column sep=small]
 \& \& \& a=(0,1)\\
 X\arrow[r]\&\underline{\Psi}={}\&0\arrow[ur]\arrow[dr]\&\\
 \&\&\&b=(1,0)
\end{tikzcd}}

les fibres \unsure{st|sont|étant}

\[
 X_0=U\cap U',\quad X_a=U\cap T',\quad X_b=U'\cap T.
\]
\deleted{\illegible{2}}

Soi\deleted{ent} $\underline U_0,\underline U'_0$ les ouverts de $\Psi$ induits par $\{0\}\times\Delta^1$, $\Delta^{\unsure{1|0}}\times\{0\}$, donc

\[
 \left\{\begin{aligned}
 \underline U_0&=\{0,a\}, &\underline U'_0&=\{0,b\}\\
 \underline U_0\cup U'_0&=\Psi, &\underline U_0^{\unsure{*|+}}\cap U'_0&=\{0\},
 \end{aligned}\right.
\]

on aura donc

\[
 \left\{\begin{aligned}
 U&=\varphi^{-1}(\underline U_0), & U'&=\varphi^{-1}(\underline U'_0)\\
 U\cap U'&=\varphi^{-1}(\underline U_0\cap U'_0)=\varphi^{-1}(\{0\})=X_{\unsure{0|e}}
 \end{aligned}\right.
\]

\scanpage{44}
% archive sheet number: 43

\mspage{90}

\textit{sexe est partout dans la psyché} (et que celle-ci le nie et ferme les yeux là où il est n'y change rien), et plus que nulle part ailleurs, il est dans son pouvoir de créer. Et d'autre part, \textit{la psyché et l'âme sont partout dans la réalité}.

Il en est ainsi, tout au moins, et de façon évidente, quand il s'agit de la réalité telle qu'elle est perçue ou conçue par la psyché, c'est-à-dire de la représentation qu'elle se fait de \margin{\handwritten{$\bullet$}}\deleted{\unsure{notr}}\handwritten{cett}e réalité, de sa relation à elle. Car l'image de la réalité que l'âme se construit aux différents niveaux conscient et inconscients est tout autant et plus une création de l'âme, qu'un simple reflet plus ou moins fidèle d'une ``réalité objective'' qui existerait indépendamment d'elle. (Réalité dont il n'est pas question ici de nier l'existence). C'est dire que \textit{l'état de maturité et les dispositions générales de l'âme déterminent autant et plus son image de la réalité, que cette réalité elle-même ne la détermine}.

Ceci vu, on conçoit qu'un état de peur de l'âme devant elle-même, et de fuite devant la réalité d'elle-même, puisse et doive entraîner une peur et une fuite devant \textit{toute} réalité. Nous pouvons dire aussi que dans notre relation à elle, \textit{toute réalité agit à la façon d'un} miroir \textit{dans lequel la psyché voit reflétée une image d'elle-même}. (Et si cette image est rarement perçue au \margin{\handwritten{$\bullet$}}niv\inserted{\handwritten{e}}au conscient, elle l'est toujours tout au moins dans l'inconscient profond, y compris dans l'inconscient moyen-profond, où se situe la peur primordiale.) En effet, cette réalité ne peut être distinguée par la psyché de l'image qu'elle s'en fait, c'est-à-dire aussi de sa relation à elle ; et nécessairement et qu'elle le veuille ou non, cette image et cette relation la reflètent.

Cette image de nous-même que nous renvoie la réalité nous implique d'une façon d'autant plus névralgique, et elle déclenche des réactions de peur et de fuite d'autant plus énergiques ou violentes, que cette réalité nous apparaît comme plus ``proche'' de nous. Le plus souvent, ces réactions sont les plus véhémentes ou les plus implacables quand la réalité concerne notre relation soit à nous-même, soit à un être proche : à un parent, \margin{\handwritten{$\bullet$}}\deleted{\unsure{a|o}}\handwritten{\unsure{a|q}}u conjoint, à \margin{\handwritten{$\bullet$}}\deleted{\unsure{votre|notre}}\handwritten{notre} enfant - et ceci d'autant plus que la relation est forte, que l'énergie investie en elle est considérable. Mais c'est surtout \textit{la force de notre relation à telle ou telle réalité, l'intensité d'intérêt et l'énergie que nous y investissons, qui sont} ici \textit{déterminantes}. Ainsi le plus souvent, des spéculations philosophique de haut vol\margin{\handwritten{$\bullet\ \bullet$}}\deleted{\unsure{,}}\handwritten{/} aur\deleted{\unsure{a}}\handwritten{\unsure{o|a}}nt beau concerner les questions les plus cruciales (et brûlantes certes !) que l'esprit humain se soit jamais posé sur lui-même, sur Dieu et sur l'âme, ces nobles envolées ne renverront à l'âme qu'une image d'elle pâle voire imperceptible, qui n'aura vertu ni de l'interpeller ni de l'effrayer. Mais même dans une science aussi abstraite et aussi éloignée en apparence de toute réalité humaine que la mathématique, pour peu que notre relation à la belle dame soit ardente et aussitôt elle devient pour nous névralgique, et l'image qu'elle nous renvoie de nous-même forte et véhémente. Et la réaction de peur devant cette image sera d'autant plus viscérale et forte que

\scanpage{45}
% archive sheet number: 44

\mspage{7}
On a aussi

\[
 U\simeq \deleted{\unsure{X}}X_{/a},\quad U'=X_{/b},\quad U\cap U'=X_{/0}.
\]

Ainsi le diagramme dans Cat

\diagram{\begin{tikzcd}[ampersand replacement=\&]
 \& U\\
 \text{(*)}\qquad U\cap U'\arrow[ur,"i_0^{\unsure{1|'}}"]\arrow[dr,"i_0"'] \&\\
 \& U'
\end{tikzcd}}

de type $\underline\Psi$ n'est-il autre que le foncteur $\underline\Psi\longrightarrow\mathrm{Cat}$ qui définit la cat \inserted{cofibrée scindée} $\Psi(X/\underline\Psi)=\overline X$ sur $\Psi$

\diagram{\begin{tikzcd}[ampersand replacement=\&]
 X\arrow[rr,"\alpha"]\arrow[dr,no head,"\varphi"']\&\&\overline X\arrow[dl,no head,"\overline\varphi"]\quad \alpha\in W_{\underline\Psi}^{\unsure{g|D|G}}\\
 \&\underline\Psi\&
\end{tikzcd}}

\[
 U=\overline X_a,\quad U'=\overline X_b,\quad U\cap U'=\overline X_0,
\]

les foncteurs $i_0^1,i_0$ de ces étant les foncteurs changement de base relatifs à $0\xrightarrow{\underline i_0^1}a$ et $0\xrightarrow{\underline i_0^2}b$. \deleted{D'autre}

D'autre part $\overline\varphi$ est \uline{propre}, donc $\overline\varphi_*$ se calcule ``fibre par fibre'', en particulier, pour coeff. const. $\xi_{\overline X}$, avec $\xi\in\mathcal A$

\[
 \overline\varphi_*(\xi_{\overline X})(s)=
 \left\{\begin{aligned}
  &H_{\underline D}^{\bullet}(U\cap U',\xi)&&\text{si }s=0\\
  &H^\bullet_{\underline D}(U\unsure{,|)}\xi)\deleted{*}&&\text{si }\xi=a\\
  &H^\bullet_{\underline D}(U',\xi)&&\text{si }s=b
 \end{aligned}\right.
\]

\scanpage{46}
% archive sheet number: 45

\mspage{8}
Donc le diagramme squelettique

$(*)$\quad\diagram{\begin{tikzcd}[ampersand replacement=\&]
H^\bullet_{\underline D}(U,\xi)\arrow[dr] \& \& H^\bullet_{\underline D}(U',\xi)\arrow[dl,no head]\\
\& H^\bullet_{\underline D}(U\cap U',\xi) \&
\end{tikzcd}},

de type $\underline\Psi^{\mathrm{opp}}$, \inserted{provient-il} d'un diagramme concret

\[
\gamma_0\deleted{\unsure{\xi|\varphi}}\overset{\mathrm{déf}}{=}\bar\varphi_*(\unsure{a|\xi}_{\bar X})\in\underline D(\underline\Psi).
\]

Considérons l'inclusion

$\underline\Psi\underset{*}{\overset{\underline k}{\subset}}\Delta^1\times\Delta^1={}$\diagram{\begin{tikzcd}[ampersand replacement=\&]
\& (0,1)\arrow[dr,no head] \&\\
(0,0)\arrow[ur,no head]\arrow[dr,no head] \& \& (1,1)\\
\& (1,0)\arrow[ur] \&
\end{tikzcd}}

déduite de $\Psi$ en lui rajoutant l'élément final manquant $(1,1)$. Soit $\underline k$ l'inclusion, et considérons le ``carré cartésien''

\[
\gamma=\underline k_*(\gamma_0)\in\underline D(\Delta^1\times\Delta^1)
\]

\deleted{Donc les fibres de $\gamma$ en les points de} On a

\[
\underline k^*(\gamma)\simeq\gamma_0
\]

donc \unsure{le|la} \deleted{fibre de $\gamma$ en} squelett. induit par $\gamma$ sur $\underline\Psi$ n'est autre que $(*)$ ci-dessus. Mais je dis que le diagramme squelettique sur $\Delta^1\times\Delta^1$ tout entier s'identifie à

\mspage{9}
$(**)$\quad\diagram{\begin{tikzcd}[ampersand replacement=\&]
\& H^\bullet_{\underline D}(\deleted{\unsure{U}}\bar X,\xi)\arrow[dl,"i^*"' ]\arrow[dr,"i'{}^*"] \&\\
H^\bullet_{\underline D}(U,\xi)\arrow[dr,"i'_0{}^*"' ] \& \& H^{\unsure{\prime|\bullet}}_{\underline D}(U',\xi)\arrow[dl,"i_0^*"]\\
\& H^\bullet_{\underline D}(U\cap U',\xi) \&
\end{tikzcd}}

% slanted
\margin{\uline{NB} On a un énoncé analogue, pour \uline{tt} $\xi\in\underline D(X)$, pas seulement $\xi$ constant\ldots}

Notons pour ceci que si \deleted{que}

\[
\psi=\underline k\circ\unsure{y|\varphi}\colon\bar X\longrightarrow\Delta^1\times\Delta^1=\underline Q
\]

\deleted{\unsure{on}} $\gamma=\underline k_*(\gamma_0)=\underline k_*(\varphi_*(\xi_{\bar X}))=\psi_*(\xi_{\bar X})$, d'où

\[
H^\bullet_{\underline D}(\Delta^1\times\Delta^1,\gamma)\simeq H^\bullet_{\underline D}(\bar X,\xi)\xrightarrow{\sim}H^*_{\underline D}(X,\xi),
\]

\margin{car $\alpha\colon X\to\bar X$ est $\in W$}

D'autre part, comme $e=(1,1)$ élément final de $\Delta^1\times\Delta^1=\underline Q$, donc $e\to\underline Q$ \unsure{coasph.|asph.},

\[
\deleted{\unsure{H}}\quad\underbrace{H^\bullet_{\underline D}(Q,\gamma)}_{H^*_{\underline D}(\bar X,\xi)}\simeq\gamma(e)=\gamma(1,1)\deleted{\unsure{\in A}},
\]

d'où l'assertion.

Ça doit faire partie, d'autre part, des propriétés générales des carrés cartésiens \deleted{que} \inserted{``dans $A$''} (comme les), que

% slanted
\margin{\uuline{à vérifier}, car à mettre \unsure{en axiomes|en revenant|comme axiome} sur les dérivateurs.}

\[
i_0^*\ \text{iso}\Longleftrightarrow i^*\ \text{iso},
\]

Or comme c'est vrai pour \uline{tt} $\xi\in A,=\underline D(e)$, cela donne

\[
i_0\in W_{\underline D}\Longrightarrow{i}\in W_{\underline D}.
\]

\scanpage{47}
% archive sheet number: 46

\mspage{10}
\deleted{(Idem pour les $D$-équivalences \deleted{fortes}\inserted{``complètes''}, i.e. des coeff. \uline{loc} constants $\xi$ au lieu de constants. Mais il faut supposer que l'on sache déjà que $X\to\bar X$ est $D$-équi. \unsure{2-complète|2-complètes},}

Pour avoir l'implication inverse

\[
i\in W_{\underline D}\Longrightarrow i_0\in W_{\underline D},
\]

il faudrait savoir que le diagramme $(**)$, plus \deleted{\unsure{point|précis}} précisément l'é\deleted{\unsure{lément}} \deleted{diagramme}\inserted{carré} substantiel $\gamma$ qui lui donne naissance, est non seulement \uline{car}tésien, mais aussi \uline{cocartésien}. ``Moralement'', qu'il soit cartésien me semble être la traduction cohomologique du fait que le \deleted{diagramme}\inserted{carré}

\diagram{\begin{tikzcd}[ampersand replacement=\&]
\& U\arrow[dl] \&\\
U\cap U' \& \& X\arrow[ul]\arrow[dl]\\
\& U'\arrow[ul] \&
\end{tikzcd}}

\unsure{dans|ds} Cat \deleted{\unsure{est}} est \uline{cocartésien} dans un sens suffisamment fort. Mais il est également cartésien, de façon triviale, et on peut s'attendre (sous certaines conditions\ldots) que les diagrammes ``transformés'' par $H^\bullet_{\underline D}(-,\xi)$ sont cocartésiens\ldots

J'ai envie de regarder d'emblée $X$ comme sur $\underline Q=\Delta^1\times\Delta^1$, plutôt que sur \unsure{Nommons|Nommant} $\underline\Psi=\underline Q-\{e\}$ de $\underline Q$, alors

\[
\underline Y_{\underline Q}(X)|\underline\Psi\simeq Y_{\underline\Psi}(X)\ \text{déjà examiné.}
\]

\scanpage{48}
% archive sheet number: 47

\mspage{11}
mais sa quatrième fibre est

\[
\bar X_{(1,1)}=X_{/(1,1)}\simeq X,
\]

donc le diagramme \inserted{dans Cat} du type $\underline Q$ donnant naissance à $Y_{\underline Q}(X)$ n'est autre que

\diagram{\begin{tikzcd}[ampersand replacement=\&]
\& U\arrow[dr,no head] \&\\
U\cap U'\arrow[ur]\arrow[dr] \& \& X\\
\& U'\arrow[ur,no head] \&
\end{tikzcd}}

Mais ce diagramme \inserted{comm. dans Cat} n'est bien sûr pas quelc, sa propriété essentielle, pour les propriétés que nous avons en vue, me semble la suivante :

\diagram{\[
\left(\underbrace{\int
\begin{tikzcd}[ampersand replacement=\&,column sep=small,row sep=small]
\& U\\
\deleted{\unsure{U\cap U'|V}}\arrow[ur]\arrow[dr,no head] \&\\
\& U'
\end{tikzcd}}_{\substack{\unsure{\parallel|\Downarrow}\,\mathrm{déf}\\Y_{\underline\Psi}(X)\ \text{en tant}\\\text{qu'objet de Cat sans plus}}}\right)\longrightarrow X
\]}

\margin{(flèche définie pour tt diagramme comm.}

\diagram{\begin{tikzcd}[ampersand replacement=\&,column sep=small,row sep=small]
\& U\arrow[dr] \&\\
V\arrow[ur]\arrow[dr,no head] \& \& X\\
\& U'\arrow[ur] \&
\end{tikzcd}}

\margin{dans Cat}

est dans $W$. Et le fait général pour y satisfaire est le suivant

``\uline{Principe}'' $(?)$ Soit

\diagram{\begin{tikzcd}[ampersand replacement=\&,row sep=small]
\& X_a\\
X_0\arrow[ur]\arrow[dr] \&\\
\& X_b
\end{tikzcd}}

un diagramme dans Cat du type $\underline\Psi$, soit $\mathcal X$ la cat. cofibrée sur $\underline\Psi$ qu'il définit. Pour tout $\xi$

\scanpage{49}
% archive sheet number: 48

\mspage{12}
$\xi\in\underline D(\mathcal X)$, considérons (si $\varphi\colon\mathcal X\longrightarrow\underline\Psi$ \unsure{est} la projection)

\[
\begin{aligned}
\gamma_0&\overset{\deleted{\unsure{k}}\ \mathrm{déf}}{=}\varphi_*(\xi)\in\underline D(\underline\Psi)\\
\gamma&=\underline k_*(\gamma_0)\in D(\underline Q);
\end{aligned}
\]

donc

\[
\begin{aligned}
\gamma(0,0)&\simeq\deleted{\unsure{X_0}}H^\bullet_{\underline D}(X_0,\xi)\\
\gamma(0,1)&\simeq H^\bullet_{\underline D}(X_a,\xi)\\
\gamma(1,0)&\simeq H^\bullet_{\underline D}(X_b,\xi)\\
\gamma(1;1)&\simeq H^\bullet_{\underline D}(\mathcal X,\xi)
\end{aligned}
\]

et \deleted{diagramme} le carré squelettique associé est

$(\alpha)$\quad\diagram{\begin{tikzcd}[ampersand replacement=\&]
\& H^\bullet_{\underline D}(\mathcal X,\xi)\arrow[dl]\arrow[dr] \&\\
H^\bullet_{\underline D}(X_a,\xi)\arrow[dr] \& \& H^\bullet_{\underline D}(X_b,\xi)\arrow[dl]\\
\& H^\bullet_{\underline D}(X_0,\xi) \&
\end{tikzcd}.}

(NB\quad\diagram{\begin{tikzcd}[ampersand replacement=\&]
\& X_a\arrow[dr] \&\\
X_0\arrow[ur]\arrow[dr] \& \& \mathcal X\\
\& X_b\arrow[ur] \&
\end{tikzcd}}

n'a rien de commun, dans $\mathcal X$\ldots{} $X_a\cap X_b=\phi$. Mais

\diagram{\begin{tikzcd}[ampersand replacement=\&]
\& X_{/a}\arrow[dr,no head] \&\\
X_{/0}\arrow[ur,no head]\arrow[dr,no head] \& \& \mathcal X\\
\& X_{/b}\arrow[ur,no head] \&
\end{tikzcd}}

est comm., et le carré $(\alpha)$ est aussi déduit de
\mspage{13}
celui-là par $H^\bullet_{\underline D}(-,\xi)$.

L'élément $\gamma\in\underline D(Q)$ est cartésien par construction, donc le ``carré'' $(\alpha)$ \unsure{aussi|assoc.}. Ça serait le cas pour \uline{tt} cat. $\mathcal X$ \inserted{presque} au dessus de $\underline\Psi$, sans avoir à le supposer cofibré. Mais le fait que $\mathcal X$ soit cofibré doit donner au carré $(\alpha)$ la vertu suppl. qu'il est aussi \uline{cocartésien}. J'ai l'impression que ça ne peut absolument pas découler de l'un des axiomes Der 1 -- Der 5,5' vus jusqu'à présent. Je vais le mettre comme

\uuline{``Der 6''} \inserted{(provisoire)} \deleted{Soit} si $\mathcal X$ sur $\underline\Psi$ est cofibré, pour $\xi$ \deleted{const. et \unsure{dérivé|dérivée}, la} \inserted{quelconque} \deleted{\unsure{constant}} dans $\underline D(\mathcal X)$, \deleted{la} l'élément $\psi_*(\xi)\in\underline D(\underline Q)$, (où $\psi=(\mathcal X\xrightarrow{\varphi}\underline\Psi\lhook\joinrel\xrightarrow{\underline k}\underset{\substack{\shortparallel\\\Delta^1\times\Delta^1}}{\underline Q})$ est non seulement cartésien, mais aussi cocartésien.

% slanted
\margin{à vérifier que ça devient \unsure{déjà|clair} pour $\xi$ seulement (resp. const.), si \unsure{$D$ dérivé|$D$ cohomol.} \unsure{directe}.}

% slanted
\margin{\uline{Axiome des cofibrations} sous-entendu : si \unsure{le dérivateur|D de domaine} est orienté \unsure{en|comme} pour la codescente, où le ``bon'' axiome \unsure{(en th.?)|(\textnormal{en théorie})} est Der 6 de la page suiv.}

\scanpage{50}
% archive sheet number: 49

\mspage{14}
``\uline{Der 6'}''\inserted{(provisoire)} Si $\mathcal X$ est fibré sur

% slanted
\margin{\deleted{Je \unsure{supprime|suppose} \illegible{1}}}

$\underline\Psi^{\mathrm{opp}}=\underline\Phi={}$\diagram{$\left(\begin{tikzcd}[ampersand replacement=\&,row sep=small]
\& a\arrow[dl]\\
1 \&\\
\& b\arrow[ul]
\end{tikzcd}\right)$}

(donc donné par, $\underline Q$-équiv. près, par un diagramme

% slanted
\margin{\uline{Axiome des fibrations}}

\diagram{\begin{tikzcd}[ampersand replacement=\&,row sep=small]
\& X_a\\
X_1\arrow[ur,no head]\arrow[dr] \&\\
\& X_b
\end{tikzcd}}

dans Cat) alors pour $\xi\in\underline D(\mathcal X)$ \deleted{constant} (voir la \unsure{co|ci|c.}) \deleted{\unsure{H}}\ $\psi_!(\gamma)\in\underline D(\underline Q)$ est non seulement cocartésien (sens évident), mais aussi cartésien.

$\underline\Phi\subset\underline Q^a={}$\diagram{\begin{tikzcd}[ampersand replacement=\&,execute at end picture={\draw (Apoint.north) to[out=178,in=80] ([xshift=-4pt]Onepoint.west) to[out=-95,in=180] (Bpoint.south) to[out=0,in=-92] ([xshift=4pt]Apoint.east) to[out=88,in=0] (Apoint.north) -- cycle;}]
\& |[alias=Apoint]|a\arrow[dl] \&\\
|[alias=Onepoint]|1 \& \& 0\arrow[ul]\arrow[dl]\\
\& |[alias=Bpoint]|b\arrow[ul] \&
\end{tikzcd}}

\uline{NB} Der 6' pour $\underline D$ \deleted{signifie} n'est autre que Der 6 pour $\underline D^\circ$. Je ne vois pas comment l'un impliquerait l'autre, en dehors d'un cas exceptionnel où il y aurait \deleted{\unsure{une dualité}}\inserted{(p.ex.)} un foncteur dualisant dans $A=\underline D(e)$, et une correspondance par dualité entre cohomologie et homologie.

\deleted{Finalement}, il semblerait que la bonne relation, qui englobe Der 6 et Der 6bis, soit la suivante :

% slanted
\margin{\uuline{la condition} \deleted{\unsure{apparue|apparente}}}

% slanted
\margin{``L'\uline{axiome} \uuline{spécifique}''}

\uline{Der\unsure{+|.}6\deleted{\unsure{bis|provisoire}}} Si $\gamma\in\underline D(\Delta^1\times\Delta^1)$, $\gamma$ est cartésien ssi il est cocartésien. \uline{De plus} (\uline{à prouver} que ça en résulte) si diagramme squelettique associé est

\scanpage{51}
% archive sheet number: 50

\mspage{15}
\diagram{\begin{tikzcd}[ampersand replacement=\&]
\& \gamma(0,1)\arrow[dl,"j_0"'] \&\\
\gamma(0,0) \& \& \gamma(1,1)\arrow[ul,"i_{1}"']\arrow[dl,"j_1"]\\
\& \gamma(1,0)\arrow[ul,"i_0"'] \&
\end{tikzcd}}

\margin{De plus, on doit imposer autant pour les dérivateurs induits $\underline D_I$ $(I\in\mathrm{Diag})$, \unsure{idem|donc} dans $\underline D(I\times\Delta^1\times\Delta^1)$.}

\[
\text{alors}\quad
\left\{\begin{aligned}i_1\ \text{iso}&\Longleftrightarrow i_0\ \text{iso},\\
j_1\ \text{iso}&\Longleftrightarrow j_0\ \text{iso}.\end{aligned}\right.
\]

Cette dernière assertion résulte de la précédente (car cart. implique $\Rightarrow$, cocartésien implique $\Leftarrow$), pourvu qu'on sache ceci :

``\uline{Lemme}'' \deleted{(associé à Der 6)} : si $\xi=\deleted{\unsure{\xi}}$

\diagram{$\left(\begin{tikzcd}[ampersand replacement=\&,row sep=small]
\& \xi_a\arrow[dl,"j_0"']\\
\xi_0 \&\\
\& \xi_b\arrow[ul,"i_0"']
\end{tikzcd}\right)$}

est un diagramme substantiel tel que $i_0$ soit un iso, alors il est isom. à l'image inverse d'un $y\in\underline D(\underline U\deleted{\unsure{\Phi}})$ (où $\underline U=(\deleted{\unsure{b}}0\longrightarrow a)\simeq\Delta^1$ est l'ouvert canonique gauche de $\underline\Phi$) par la rétraction \inserted{$p$} de $\underline\Psi$ sur $\underline U$.

% margin note written sideways
\margin{\deleted{\unsure{cet axiome est lié à|cet axiome est le}}}

Il suffit de prouver que

\[
\deleted{\unsure{\xi}}\ p^*p_*(\xi)\longrightarrow\xi\quad\text{, adj.}
\]

est iso. Comme $p$ est bien sûr lisse, comme induit par $\mathrm{pr}_1\colon\Delta^1\times\Delta^1\longrightarrow\Delta^1$, sur un ouvert, mais aussi \deleted{propre} $p_*(\xi)$ se calcule fibre par fibre. On utilise que \deleted{\unsure{dans|$\leftarrow$}}

\[
A\simeq\underline D(e)\longrightarrow\underline D^{\unsure{\ell c|ec|cc}}(\Delta^1)
\]

est une équiv. de cat. (\uline{forme de Der 3})$^*$.

\scanpage{52}
% archive sheet number: 51

\mspage{16}
Dans Hot, par contre on a ceci :

Tout carré cocartésien est cartésien. Pour qu'un carré cartésien

\diagram{\begin{tikzcd}[ampersand replacement=\&]
\& \xi_0\arrow[dl,no head]\arrow[dr,no head] \&\\
\xi_a\arrow[dr,no head] \& \& \xi_b\arrow[dl,no head]\\
\& \xi_1 \&
\end{tikzcd}}

soit aussi cocartésien, il f. et s. que l'on ait

\[
\pi_0(\xi_1)=\mathrm{Im}\,\pi_0(\xi_a)\cup\mathrm{Im}\,\pi_0(\xi_b)
\]

% slanted
\margin{complètement faux. Si p.ex. $\xi_1=e$, le diagramme ci-contre : le fait homotopiquement cartésien et cocartésien pourvu que $\xi_a\simeq e$ ou $\xi_b\simeq e$}

Der 6 ou Der 6' au choix donnent W7. Quant à W8, il semble tributaire d'un principe général \inserted{(``homom. des cinq'')}

``\uline{Der 8}'' Soit un hom. d'un carré \deleted{\unsure{bi}}cartésien dans un autre, dans $A^{\prime\prime}(\Delta^1\times\Delta^1)$. Considérons

\diagram{\begin{tikzcd}[ampersand replacement=\&,column sep=small,row sep=small]
\& \xi_a\arrow[dr,no head]\arrow[ddd,"f_a"'{pos=.22}] \& \&\\
\xi_0\arrow[ur]\arrow[dr]\arrow[ddd,"f_0"'] \& \& \xi_1\arrow[ddd,"f_1"] \&\\
\& \xi_b\arrow[ur]\arrow[ddd,"f_b"] \& \&\\
\& \eta_a\arrow[dr] \& \&\\
\eta_0\arrow[ur]\arrow[dr] \& \& \eta_1 \&\\
\& \eta_b\arrow[ur] \& \&
\end{tikzcd}}

les conditions

a) $f_0$ est iso

b) $f_1$ est iso

c) $f_a$ \uline{et} $f_b$ est iso.

Alors \inserted{si les carrés sont ``isom.'' cartésiens,} \deleted{\unsure{chacune des deux}} $b+c\Rightarrow a$. S'ils sont bien cocartésiens, $a+c\Rightarrow b$. \deleted{\unsure{Si les carrés sont cartésiens}} S'ils sont bicartésiens, on a \deleted{les formules} \unsure{(il est)} de plus $a+b\Rightarrow c$ (?).

\margin{\inserted{vérifier si c'est raisonnable, pas évident}}

% slanted
\margin{Il suffit d'exiger que $a)+b)\Rightarrow c$ i.e. si $f_0,f_1$ iso, alors $f_a$ et $f_b$ iso. \unsure{Ceci en|Tout ceci en} évidence.}

\uuline{NB} \deleted{La \unsure{form.} $a+c\Rightarrow b$, \unsure{est|étant} cocartésien donc $b+c\Rightarrow a$ s'il suffit cartésien et c'est évident. Le cas non évident est le cas $a+b\Rightarrow c$}

% slanted
\margin{version \unsure{renforcée} de W7 page 4.}

\mspage{17}
\uline{W7bis} Soit

\diagram{\begin{tikzcd}[ampersand replacement=\&,row sep=small]
\& X_a\\
X_0\arrow[ur]\arrow[dr] \&\\
\& X_b
\end{tikzcd}}

diagramme dans Cat, soit

$\underline\Phi={}$\diagram{\begin{tikzcd}[ampersand replacement=\&,row sep=small]
\& a\arrow[dl]\\
0 \&\\
\& b\arrow[ul]
\end{tikzcd}}

et $\mathcal X$ la catégorie fibrée sur $\underline\Phi$ définie par le diagramme. On a un diagramme commutatif dans Hot$_W$ :

\diagram{\begin{tikzcd}[ampersand replacement=\&]
\& X_0\arrow[dl,"i_a"']\arrow[dr,"i_b"] \&\\
X_a\arrow[dr,"i_{b'}"'] \& \& X_b\arrow[dl,"i_{a'}"]\\
\& \mathcal X \&
\end{tikzcd}}

% slanted
\margin{\uline{NB} Autre présentation. On \uline{part} de $\mathcal X$, fibrée sur}

\diagram{\begin{tikzcd}[ampersand replacement=\&,row sep=tiny,column sep=tiny]
\& a\arrow[dl]\\
0 \&\\
\& b\arrow[ul]
\end{tikzcd}}

(NB Dans Cat, \deleted{je présume que} les deux flèches composées sont homotopes, par un chemin de longueur 2 \deleted{un} elles sont homotopes tous deux, par une homotopie de long. 1, à ``l'inclusion'' $X_0\lhook\joinrel\longrightarrow\mathcal X$)

Alors \deleted{\unsure{iso|tous}}

\[
\begin{aligned}
i_a\in W&\Longleftrightarrow i_{a'}\in W\\
i_b\in W&\Longleftrightarrow i_{b'}\in W
\end{aligned}
\]

\uline{W8bis} Soit un diagramme comm. dans Cat

\diagram{\begin{tikzcd}[ampersand replacement=\&,column sep=small,row sep=small]
\& X_a\arrow[dr,dashed]\arrow[ddd,"f_a"'{pos=.22}] \&\\
X_0\arrow[ur]\arrow[dr]\arrow[ddd,"f_0"'] \& \& \mathcal X\arrow[ddd,"f"]\\
\& X_b\arrow[ur,dashed]\arrow[ddd,"f_b"] \&\\
\& X'_a\arrow[dr,dashed] \&\\
X'_0\arrow[ur]\arrow[dr] \& \& \mathcal X'\\
\& X'_b\arrow[ur,dashed] \&
\end{tikzcd}}

\margin{(traits pleins)}

% slanted
\margin{équivalent : l'énoncé \unsure{direct}, avec $\mathcal X,\mathcal X'$ cofibrées sur $\underline\Psi$ et un \uline{cocartésien} $\mathcal X\to\mathcal X'$.}

% slanted
\margin{à examiner, ici et dans W7bis, si l'hyp. ``$\unsure{f'|f}$ cartésien'' (resp. cocartésien) est nécessaire ou superflue ?}

On le complète par un diagramme commutatif dans Hot$_W$ (``par construction'' de $\mathcal X,\mathcal X'$\ \deleted{dans}, où pour

\diagram{$X=\displaystyle\int\left(\begin{tikzcd}[ampersand replacement=\&,row sep=small,column sep=small]
\& X_a\\ X_0\arrow[ur,no head]\arrow[dr] \&\\ \& X_b
\end{tikzcd}\right)$}

). Considérons les cond.

\[
\left\{\begin{aligned}
\text{a)}&\quad f_{\unsure{*|1}}\in W\\
\text{b)}&\quad f_0\in W\\
\text{c)}&\quad f_a\ \text{et}\ f_b\in W
\end{aligned}\right.
\]

Alors deux \deleted{\unsure{deux|des}} ces conditions impliquent la troisième.

\scanpage{53}
% archive sheet number: 52

\mspage{18}
\deleted{Version forte}\inserted{Corollaire de W7bis} (\deleted{\unsure{vers.|très|sous}} forme plus forte en apparence)

\uline{W7\unsure{ter|bis}} Soit $\mathcal X$ lisse sur

$\Phi={}$\diagram{\begin{tikzcd}[ampersand replacement=\&,column sep=small,row sep=small]
\& a\arrow[dl]\\
0\arrow[dr] \&\\
\& b
\end{tikzcd}}

, d'où diagramme (comm. dans Hot$_W$)

\diagram{\begin{tikzcd}[ampersand replacement=\&]
\& X \&\\
X_a\arrow[ur,"i_{b'}"'] \& \& X_b\arrow[ul,"i_{a'}"]\\
\& X_0\arrow[ul,"i_a"']\arrow[ur,"i_b"'] \&
\end{tikzcd}}

Alors

\[
\begin{aligned}
i_a\in W&\Longleftrightarrow i_{a'}\in W\\
(\text{donc}\quad i_b\in W&\Longleftrightarrow i_{b'}\in W)
\end{aligned}
\]

(NB Cela implique l'axiome dual W7'ter, avec $X$ propre sur

$\Psi={}$\diagram{\begin{tikzcd}[ampersand replacement=\&,column sep=small,row sep=small]
\& a\arrow[dl]\\
0 \&\\
\& b\arrow[ul]
\end{tikzcd}} )

Corollaire de W8bis (forme renforcée en apparence).

\uline{W8ter} Soient $\mathcal X,\mathcal X'$ lisses sur

$\underline\Phi={}$\diagram{\begin{tikzcd}[ampersand replacement=\&,column sep=small,row sep=small]
\& a\arrow[dl]\\
0 \&\\
\& b\arrow[ul]
\end{tikzcd}}

, et $f\colon\mathcal X\longrightarrow\mathcal X'$ un $\underline\Phi$-morphisme. Soient $f_a,f_b,f_0$ induit sur les fibres, et les conditions

\[
\begin{aligned}
\text{a)}\ & f\in W\\
\text{b)}\ & f_0\in W\\
\text{c)}\ & f_a\ \text{et}\ f_b\in W.
\end{aligned}
\]

Alors deux de ces conditions impliquent la troisième.
